\subsection{推理异步化：本章重点}\label{sec:async-inference}

\subsubsection{设计动机：让推理与执行真正并行}\label{sec:async-inf-motivation}

回到 4.1.2 节指出的现象：在同步流水线下，每个``推理 + 执行 chunk''循环里
有 3.3 s 的 CPU 闲置窗口。如果能把下一帧的推理重叠到当前 chunk 的执行期里，
吞吐有望从 ``$T_{\text{infer}} + N/F$''---限制（约 19 FPS）
恢复到仅受执行频率约束的 30 FPS。这就是推理异步化的核心设计动机。

与 4.2 / 4.3 节讨论的舵机异步化不同的是：推理过程纯粹是 Python+PyTorch
计算，不与任何半双工硬件资源耦合，理论上具备真正并行的物理基础；
LeRobot 上游也已为此提供了完整的双进程实现，无需自行从零开发。
因此本节的工作不再是``评估能不能做''，而是``评估在树莓派 5 + 当前 ACT 模型
的具体参数下，做了能不能带来实际收益''。

\subsubsection{LeRobot 官方双进程架构}\label{sec:async-inf-arch}

LeRobot 把推理异步化拆成两个独立进程，通过 gRPC 通信：

\begin{itemize}[leftmargin=2em]
  \item \texttt{robot\_client.py}：主进程，连接硬件，跑控制循环，
        采观测发给 server，从 server 收回 \texttt{action\_chunk} 后
        按帧发给舵机执行；
  \item \texttt{policy\_server.py}：推理进程，加载 ACT 模型权重，
        收到观测后运行 \texttt{policy.predict\_action\_chunk()}，
        把动作序列序列化（pickle）后回传。
\end{itemize}

两进程之间的关键数据载体是 \texttt{TimedAction} 与 \texttt{TimedObservation}
（见 \texttt{helpers.py}），都包含 \texttt{timestamp} 和 \texttt{timestep} 字段，
分别用于跨进程网络延迟测量和帧号对齐。

在 \texttt{robot\_client} 内部，主线程跑控制循环、后台线程
（\texttt{receive\_actions}）阻塞等待 server 回包；两个线程通过
带锁的 \texttt{action\_queue}（\texttt{queue.Queue}）共享数据。
关键状态变量包括：
\begin{itemize}[leftmargin=2em]
  \item \texttt{latest\_action}：已发给舵机的最新帧号，
        初始为 -1；
  \item \texttt{action\_chunk\_size}：server 每次返回的 chunk 长度，
        首次收到 chunk 后更新；
  \item \texttt{must\_go}：\texttt{threading.Event}，
        队列即将耗尽时强制 server 立即推理。
\end{itemize}

\subsubsection{三个核心机制}\label{sec:async-inf-mechanisms}

\textbf{机制一——\texttt{chunk\_size\_threshold}（提前推理）}：
client 主循环每帧检查 \texttt{action\_queue.qsize() / action\_chunk\_size}
是否已经 $\leq$ \texttt{chunk\_size\_threshold}（默认 0.5）。
即\textbf{队列剩 50\% 时就开始采下一帧观测、向 server 发推理请求}，
而不是等队列彻底空了才开始。设计意图是让``新 chunk 到货的时刻''
正好接上``旧 chunk 耗尽的时刻''，避免主循环 stall。

\textbf{机制二——Temporal Ensemble（多 chunk 加权融合）}：
新 chunk 到达时不是直接覆盖队列，而是按 \texttt{timestep} 对齐。
对于``新旧 chunk 都覆盖到的同一未来帧''，使用用户配置的
\texttt{aggregate\_fn}（默认取新值，可改为加权平均）做融合；
对于``新 chunk 覆盖到、但已经被主线程执行过''的帧（\texttt{timestep} $\leq$
\texttt{latest\_action}）则直接丢弃；对于``纯未来帧''则入队。
这是 ACT 原论文提出的 temporal ensemble 在异步场景下的自然实现，
用于减少 chunk 切换时的动作抖动（jitter）。

\textbf{机制三——\texttt{must\_go}（队列空兜底）}：
如果推理慢到队列真的空了，client 在采观测时把 \texttt{must\_go=True}
随观测一起发给 server；server 端跳过所有去重检查，强制把这次观测放入推理队列；
client 收到新 chunk 后再 \texttt{must\_go.set()} 重置。
这条路径属于异常兜底，正常稳态下不应频繁触发——
\texttt{must\_go} 触发频次因此是判断流水线健康度的重要指标
（4.5 节实验会重点采集）。

\subsubsection{可行性不等式 \boldmath$N \geq 2 T_{\text{infer}} F$\unboldmath\ 的推导}\label{sec:async-inf-inequality}

异步推理的稳态行为受三个参数控制：chunk\_size $N$、
单次推理耗时 $T_{\text{infer}}$、目标执行频率 $F$。
本节从单次推理产出的``有效帧数''出发推导稳态条件。

设主线程在第 0 帧完成对当前 chunk 的耗尽前 50\%（即触发阈值），
开始采观测发给 server。server 推理耗时 $T_{\text{infer}}$，
这段时间内主线程消耗 $T_{\text{infer}} \cdot F$ 帧。当新 chunk 到货时，
chunk 内的前 $T_{\text{infer}} \cdot F$ 帧已经过期（其 \texttt{timestep}
$\leq$ \texttt{latest\_action}），\textbf{在 client 端被直接丢弃}；
真正能进入队列、被未来执行的有效帧数为：
\begin{equation}
  N_{\text{eff}} = N - T_{\text{infer}} \cdot F.
\end{equation}

下一次推理周期需要主线程消耗另外 $T_{\text{infer}} \cdot F$ 帧才能再次到货，
因此\textbf{稳态不 stall 的条件}是``一次推理产出的有效帧数 $\geq$
下一次推理期间消耗的帧数''：
\begin{equation}
  N - T_{\text{infer}} \cdot F \;\geq\; T_{\text{infer}} \cdot F
  \quad\Longleftrightarrow\quad
  \boxed{\,N \;\geq\; 2 \, T_{\text{infer}} \, F\,}.
\end{equation}

为直观展示该不等式的物理含义，图~\ref{fig:async-inf-three-cases} 给出
$N=100$、$F=30$ 时三种代表性 $T_{\text{infer}}$ 取值下的稳态时间线
（一格 $=$ 10 帧；\texttt{████} 表示正常执行，\texttt{░░░░} 表示 stall）。

\begin{figure}[htbp]
  \centering
  \begin{minipage}{0.95\linewidth}
    \scriptsize\ttfamily
    \textbf{情况 A: T\_infer = 30 帧 (1 s)\quad 不等式满足，稳态 30 FPS}\\[2pt]
    执行: ████ ████ ████ ████ ████ ████ ████ ████ ████ ████ ████ ████ ████\\
    推理:\hspace{8em}P1（30帧）\hspace{6.5em}P2（30帧）\\[6pt]

    \textbf{情况 B: T\_infer = 50 帧 (1.67 s)\quad 临界点，无 ensemble 余量}\\[2pt]
    执行: ████ ████ ████ ████ ████ ████ ████ ████ ████ ████ ████ ████ ████\\
    推理:\hspace{8em}P1（50帧）\hspace{8.5em}P2（50帧）\\[6pt]

    \textbf{情况 C: T\_infer = 70 帧 (2.33 s)\quad 不等式不满足，周期性 stall}\\[2pt]
    执行: ████ ████ ████ ████ ████ ████ ████ ████ ████ ████ ░░░░ ░░░░ ████\\
    \hspace{2em}└── chunk A 100 帧 ──┘\hspace{1em}└stall 20 帧┘\hspace{0.5em}└── chunk B 50 帧 ──...
  \end{minipage}
  \caption{$N=100$、$F=30$ 时三种 $T_{\text{infer}}$ 下的稳态时间线}
  \label{fig:async-inf-three-cases}
\end{figure}

值得指出的是情况 C 下的一个微妙之处：stall 期间机器人停止执行，
\texttt{latest\_action} 被冻结，因此当新 chunk 到货时，
``过期帧''的判定基准并不是``物理时刻已经走过的帧数'' $T_{\text{infer}} \cdot F$，
而是 stall 之前最后一次执行到的帧号。
也就是说，stall 的那段时间``反过来保护''了新 chunk 的有效帧不被丢弃。
具体到情况 C，新 chunk 时间戳 $[50, 149]$ 到货时
\texttt{latest\_action} = 99（chunk A 最末帧），
有效帧数 = $|[100, 149]|$ = 50 而非 30。
这一点必须在公式推导中保留，否则时间线与公式会出现不自洽。

\subsubsection{不等式在本平台的代入}\label{sec:async-inf-substitute}

代入第三章的实测参数 $T_{\text{infer}}$ = 2 s、$F$ = 30 FPS：
\begin{equation}
  2 T_{\text{infer}} F \;=\; 2 \times 2 \times 30 \;=\; 120 \;>\; 100 \;=\; N.
\end{equation}

不等式不满足。理论稳态 FPS 退化为：
\begin{equation}
  F_{\text{actual}} \;\approx\; \frac{N - T_{\text{infer}} F}{T_{\text{infer}}}
  \;=\; \frac{100 - 60}{2} \;=\; 20\ \text{FPS},
\end{equation}
与同步模式的实测约 19 FPS 几乎持平。也就是说在当前参数下，
异步推理理论上几乎不会带来吞吐提升，反而会引入双进程、gRPC 通信、
锁竞争和 CPU 争抢等额外复杂度。这一理论预测正是 4.5 节实验要验证的对象。

\subsubsection{树莓派 5 单机部署的 CPU 争抢}\label{sec:async-inf-cpu-contention}

异步推理在理论分析之外还有一项实操层面的注意点：
树莓派 5 只有 4 个 ARM Cortex-A76 核心。
ACT 推理通过 PyTorch 的 OpenMP 后端默认会吃满所有核心，
而控制循环、串口 IO、相机解码也需要约 1 核。如果 \texttt{policy\_server}
和 \texttt{robot\_client} 同机部署，两者会在核心层面互相挤压，
推理线程被频繁让出 CPU 导致 $T_{\text{infer}}$ 抬高。

为隔离两进程对 CPU 的争抢，4.5 节实验将采用以下控制措施：
\begin{itemize}[leftmargin=2em]
  \item \texttt{taskset -c 0-2} 把 \texttt{policy\_server} 钉在核心 0--2；
  \item \texttt{taskset -c 3} 把 \texttt{robot\_client} 钉在核心 3；
  \item 在 \texttt{policy\_server} 内显式 \texttt{torch.set\_num\_threads(3)};
  \item 在 \texttt{robot\_client} 内 \texttt{cv2.setNumThreads(1)}
        避免 OpenCV 内部多线程与控制循环抢 CPU。
\end{itemize}

需要承认，3 核推理相对 4 核独占大约慢 25\%，
$T_{\text{infer}}$ 会相应抬高，进一步压缩不等式的理论余量。
这是单机部署的固有代价，也是 4.5 节实验中需要单独测量的量。
